% !TeX root = ../教材题库_全部条目_转写工作稿.tex
% Source: XJ-B1-C04-EX-P134-Y0942
% 第4章 / 4.4 函数与方程 / 教材 p.134、135 / PDF 第141、142页
% OCR未校对：公式与文字均可能有误；未完成数学TeX转写；禁止直接发给学生
% 原题及相邻原解截图（核对用）：../题目原图/XJ-B1-C04-EX-P134-Y0942-01.webp, ../题目原图/XJ-B1-C04-EX-P134-Y0942-02.webp
% 本文件只在转写完成后作为题干使用；OCR 条目尚不是可靠题干。
\begin{bankitem}{4.4 例题 例题5}{第4章 / 4.4 函数与方程 / 教材 p.134、135 / PDF 第141、142页}{XJ-B1-C04-EX-P134-Y0942}{OCR未校对：公式与文字均可能有误；未完成数学TeX转写；禁止直接发给学生}
\begin{ocrdraft}
[教材p.134 / PDF第141页]
\par
0二0]求曲线y=ln和直线3，一亠r十2的交点的横坐标（误差不超过0．01）．
\par
解曲线3，一lnx和直线3，一一十2的交点的横坐标应满足等式一一十
\par
2，即lnx十一2一0，则交点的横坐标是函数，一不一ln十一2的零点．
\par
由不1)一一1<0和不2）一ln2>0可知，丆0）在区间（1，2）内有一个零点；由
\par
不单调递增可知，它只有这一个零点．用二分法计算，借助计算器可列表如下：
\par
零点所在
\par
是否符合
\par
a+b
\par
次数区间左端点a
\par
区间右端点方
\par
不m）的近似值
\par
2
\par
区间长度
\par
要求
\par
1
\par
1
\par
2
\par
一0．09
\par
否
\par
2
\par
2
\par
1.75
\par
0．31
\par
0．25
\par
否
\par
3
\par
1.75
\par
l.625
\par
0．11
\par
0．125
\par
否
\par
4
\par
1.625
\par
1.5625
\par
0．009
\par
0．0625
\par
否
\par
5
\par
1.5625
\par
l.53125
\par
一0．0
\par
0．03125
\par
否
\par
6
\par
1.53125
\par
1.5625
\par
1.546875
\par
一0．02
\par
0．0巧625
\par
否
\par
7
\par
1.546875
\par
1.5625
\par
1.5546875
\par
一0．004
\par
0．0078125
\par
是
\par
由上表可知，计算到第7次，由于f(l.
\par
5546875）\ensuremath{\cdot}不1.5625）<0，可知包含
\par
[教材p.135 / PDF第142页]
\par
零点的区间为匚1.5546875，1.5625]，零点所在区间长度小于（).01.于是，可取
\par
I.555作为零点的近似值．因此，交点的横坐标约为1.555（可参看图4，4一6）．
\par
y=lnX+X-2
\par
2
\par
y=lnx
\par
2
\par
3
\par
4
\par
图4．4一6
\end{ocrdraft}
\end{bankitem}
